\section{Análise de Resultados}

% Secção de gráficos e todos os cálculos, incluindo a análise dos erros, e resultado final com número correto de algarismos significativos
\subsection{Montagem 1}

\begin{table}[ht!]
	\centering
	\caption{Resultados da Montagem 1.}
	\begin{tabular}{|c|c|c|c|c|}
		\hline
		\textbf{Medição} & \textbf{Valor experimental (mA)} & \textbf{Valor teórico (mA)} & \textbf{Desvio (mA)} & \textbf{Erro (\%)} \\
		\hline
		$I_{AB}$         & $44,70 \pm 0.01$                 & $45,00$                     & $-0,30$              & $0,67$             \\
		\hline
		$I_{BC}$         & $15,06 \pm 0.01$                 & 15,00                       & $+0,06$              & $0,40$              \\
		\hline
		$I_{BD}$         & $30,00 \pm 0.01$                 & 30,00                       & $ \pm 0,00$          & $0,00$             \\
		\hline
	\end{tabular}
\end{table}

\subsubsection{Cálculos:}
\[
\begin{aligned}
V_{CE}&=V_{DE}=V_0=15\,\mathrm{V} \\
I_{BC}&=\frac{V_{CE}}{R_1}=\frac{15}{1\,\mathrm{k}\Omega}=0.015\,\mathrm{A}=15\,\mathrm{mA} \\
I_{BD}&=\frac{V_{DE}}{R_2}=\frac{15}{500\,\Omega}=0.030\,\mathrm{A}=30\,\mathrm{mA} \\
I_{AB}&=I_{BC}+I_{BD}=15\,\mathrm{mA}+30\,\mathrm{mA}=45\,\mathrm{mA}
\end{aligned}
\]

\[
\begin{aligned}
d_{AB}&=I_{AB,\mathrm{exp}}-I_{AB,\mathrm{teo}}
=44.70-45.00=-0.30\,\mathrm{mA}
\qquad
Er_{AB}(\%)=\frac{|d_{AB}|}{I_{AB,\mathrm{teo}}}\times100
=\frac{|-0.30|}{45.00}\times100=0.67\% \\[6pt]
d_{BC}&=I_{BC,\mathrm{exp}}-I_{BC,\mathrm{teo}}
=15.06-15.00=+0.06\,\mathrm{mA}
\qquad
Er_{BC}(\%)=\frac{|d_{BC}|}{I_{BC,\mathrm{teo}}}\times100
=\frac{|0.06|}{15.00}\times100=0.40\% \\[6pt]
d_{BD}&=I_{BD,\mathrm{exp}}-I_{BD,\mathrm{teo}}
=30.00-30.00=0.00\,\mathrm{mA}
\qquad
Er_{BD}(\%)=\frac{|d_{BD}|}{I_{BD,\mathrm{teo}}}\times100
=\frac{|0.00|}{30.00}\times100=0.00\%
\end{aligned}
\]
\subsubsection{Verificação da lei dos nós:}

De acordo com a secção \ref{sec:fundam_teorica}, considerando o nó $B$ (ponto de junção dos ramos $AB$, $BC$ e $BD$) e adotando a convenção de sinais em que a corrente entra no nó pelo ramo $AB$ e sai pelos ramos $BC$ e $BD$, tem-se:

\begin{equation*}
	\sum I = 0
\end{equation*}

ou seja,

\begin{equation*}
	I_{AB}-I_{BC}-I_{BD}=0
\end{equation*}

Recorrendo aos valores experimentais da Tabela 3:

\begin{equation*}
	(44,70 \pm 0,01) - (15,06 \pm 0,01) - (30,00 \pm 0,01)
	= (-0,36 \pm 0,03)\ \text{mA}
\end{equation*}

Visto que o valor previsto pela lei dos nós $(0,00\ \text{mA})$ não se encontra dentro do intervalo determinado experimentalmente $(-0,36 \pm 0,02)\ \text{mA}$, conclui-se que a lei dos nós \textbf{não foi} verificada dentro da incerteza experimental associada às medições.


\subsection{Montagem 2-A}
\begin{table}[H]
	\centering
	\caption{Resultados da Montagem 2-A.}
	\begin{tabular}{|c|c|c|c|c|}
		\hline
		\textbf{Medição} & \textbf{Valor experimental} & \textbf{Valor teórico} & \textbf{Desvio} & \textbf{Erro (\%)} \\
		\hline
		$I$              & $29,4 \pm 0.01$ mA          & $10,00$ mA             & $+19,4$ mA      & $194$              \\
		\hline
		$V_0$            & $15,18 \pm 0.01$ V          & 15,00 V                & $+0,18$ V       & $1,20$             \\
		\hline
		$V_1$            & $10,10 \pm 0.01$ V          & 10,00 V                & $ +0,10$ V      & $1,00$             \\
		\hline
		$V_2$            & $5,06 \pm 0.01$ V           & 5,00 V                 & $ +0,06$ V      & $1,20$             \\
		\hline
	\end{tabular}
\end{table}
\subsubsection{Cálculos:}

\[
\begin{aligned}
R_{\mathrm{eq}}&=R_1+R_2=1,0\,\mathrm{k}\Omega+500\,\Omega=1,5\,\mathrm{k}\Omega \\
I&=\frac{V}{R_{\mathrm{eq}}}=\frac{15}{1,5\times10^{3}}=1,0\times10^{-2}\,\mathrm{A}=10\,\mathrm{mA} \\
V_1&=IR_1=(10\,\mathrm{mA})(1,0\,\mathrm{k}\Omega)=10\,\mathrm{V} \\
V_2&=IR_2=(10\,\mathrm{mA})(500\,\Omega)=5,0\,\mathrm{V}
\end{aligned}
\]

\[
\begin{aligned}
d_I&=I_{\mathrm{exp}}-I_{\mathrm{teo}}
=29,4-10,0=+19,4\,\mathrm{mA}
\qquad
Er_I(\%)=\frac{|d_I|}{I_{\mathrm{teo}}}\times100
=\frac{19,4}{10,0}\times100=194\% \\[6pt]
d_{V_0}&=V_{0,\mathrm{exp}}-V_{0,\mathrm{teo}}
=15,18-15,00=+0,18\,\mathrm{V}
\qquad
Er_{V_0}(\%)=\frac{|d_{V_0}|}{V_{0,\mathrm{teo}}}\times100
=\frac{0.18}{15,00}\times100=1,20\% \\[6pt]
d_{V_1}&=V_{1,\mathrm{exp}}-V_{1,\mathrm{teo}}
=10,10-10,00=+0,10\,\mathrm{V}
\qquad
Er_{V_1}(\%)=\frac{|d_{V_1}|}{V_{1,\mathrm{teo}}}\times100
=\frac{0,10}{10,00}\times100=1,00\% \\[6pt]
d_{V_2}&=V_{2,\mathrm{exp}}-V_{2,\mathrm{teo}}
=5,06-5,00=+0,06\,\mathrm{V}
\qquad
Er_{V_2}(\%)=\frac{|d_{V_2}|}{V_{2,\mathrm{teo}}}\times100
=\frac{0,06}{5,00}\times100=1,20\%
\end{aligned}
\]



\subsubsection{Verificação da lei das malhas:}
De acordo com a secção \ref{sec:fundam_teorica}, considerando a malha definida pelo circuito, temos que:
\begin{equation*}
	\sum V = 0
\end{equation*}

Recorrendo aos valores da Tabela 1:
\begin{equation*}
	(15,18\pm0,01) - (10,10\pm0,01) - (5,06\pm0,01) = (0,02\pm0,03)
\end{equation*}

Visto que o valor experimental $(0,00\ \text{V})$ está inserido na gama de valores determinados pelas medições e pela incerteza experimental associada, conclui-se que a lei das malhas de Kirchhoff \textbf{foi} verificada.

\subsection{Montagem 2-B}

\begin{table}[H]
	\centering
	\caption{Resultados da Montagem 2-B Analise}
	\begin{tabular}{|c|c|c|c|c|}
		\hline
		\textbf{Medição} & \textbf{Valor experimental}   & \textbf{Valor teórico} & \textbf{Desvio} & \textbf{Erro \%} \\
		\hline
		$I_0 = I_2$            & $ 13,83\pm 0.01 \, \text{mA}$ & $ 13,8 \, \text{mA}$  & $+0,03$\,mA     & $0,22$           \\
		\hline
		$I_1$            & $ 2,29\pm 0.01\, \text{mA}$  & $ 2,31 \, \text{mA}$  & $-0,02$\,mA     & $0,87$           \\
		\hline
		$I_3$            & $ 11,53\pm 0.01\, \text{mA}$  & $ 11,5 \, \text{mA}$  & $+0,03$\,mA     & $0,26$           \\
		\hline
		$V_0$            & $15,17 \pm 0.01\, \text{V}$   & $ 15,00 \, \text{V}$   & $+0,17$\,V      & $1,13$           \\
		\hline
		$V_1 = V_3 $            & $1,169\pm 0.01\, \text{V}$    & $ 1,15 \text{V}$    & $-0,331$\,V      & $22,07$           \\
		\hline
		$V_2$            & $13,95\pm 0.01\, \text{V}$    & $ 13,8 \, \text{V}$    & $+0,15$\,V      & $1,09 $          \\
		\hline
	\end{tabular}
\end{table}

\subsubsection{Cálculos:}
\[
\begin{aligned}
R_{13}&=\left(\frac{1}{R_1}+\frac{1}{R_3}\right)^{-1}
=\left(\frac{1}{500}+\frac{1}{100}\right)^{-1}
=8,33\times10^{1}\,\Omega=83,3\,\Omega \\
R_{\mathrm{eq}}&=R_{13}+R_2=83,3+1,00\times10^{3}
=1,083\times10^{3}\,\Omega \\
I_0&=\frac{V_0}{R_{\mathrm{eq}}}
=\frac{15}{1,083\times10^{3}}
=1,38\times10^{-2}\,\mathrm{A}=13,8\,\mathrm{mA}=I_2 \\
V_2&=I_2R_2=(1,38\times10^{-2})(1,00\times10^{3})
=13,8\,\mathrm{V}, \qquad
V_1=V_3=V_0-V_2=15,0-13,8=1,15\,\mathrm{V} \\
I_1&=\frac{V_1}{R_1}=\frac{1,15}{500}
=2,31\times10^{-3}\,\mathrm{A}=2,31\,\mathrm{mA}, \qquad
I_3=\frac{V_3}{R_3}=\frac{1,15}{100}
=1,15\times10^{-2}\,\mathrm{A}=11,5\,\mathrm{mA}
\end{aligned}
\]

\[
d_{I_2} = I_{2,\text{exp}} - I_{2,\text{teo}}
= 13,83 - 13,8
= 0,03\ \text{mA}
\qquad
Er_{I_2}(\%) =
\frac{|d_{I_2}|}{I_{2,\text{teo}}}\times100
=
\frac{0,03}{13,8}\times100
= 0,22\%
\]

\[
d_{I_1} = I_{1,\text{exp}} - I_{1,\text{teo}}
= 2,29 - 2,31
= -0,02\ \text{mA}
\qquad
Er_{I_1}(\%) =
\frac{|d_{I_1}|}{I_{1,\text{teo}}}\times100
=
\frac{0,02}{2,31}\times100
= 0,87\%
\]

\[
d_{I_3} = I_{3,\text{exp}} - I_{3,\text{teo}}
= 11,53 - 11,5
= 0,03\ \text{mA}
\qquad
Er_{I_3}(\%) =
\frac{|d_{I_3}|}{I_{3,\text{teo}}}\times100
=
\frac{0,03}{11,5}\times100
= 0,26\%
\]

\[
d_{V_0} = V_{0,\text{exp}} - V_{0,\text{teo}}
= 15,17 - 15,00
= 0,17\ \text{V}
\qquad
Er_{V_0}(\%) =
\frac{|d_{V_0}|}{V_{0,\text{teo}}}\times100
=
\frac{0,17}{15,00}\times100
= 1,13\%
\]

\[
d_{V_1} =  V_{1,\text{exp}} - V_{1,\text{teo}}
= 1,169 - 1,5
= -0,331\ \text{V}
\qquad
Er_{V_1}(\%) =
\frac{|d_{V_1}|}{V_{1,\text{teo}}}\times100
=
\frac{0,331}{1,5}\times100
= 22,07\%
\]

\[
d_{V_2} = V_{2,\text{exp}} - V_{2,\text{teo}}
= 13,95 - 13,8
= 0,15\ \text{V}
\qquad
Er_{V_2}(\%) =
\frac{|d_{V_2}|}{V_{2,\text{teo}}}\times100
=
\frac{0,15}{13,8}\times100
= 1,09\%
\]




\subsubsection{Verificação da lei das malhas:}
De acordo com a secção \ref{sec:fundam_teorica}, considerando a malha definida pelo circuito constituído pela fonte e as resistências $R_1$ e $R_2$, temos que:

\begin{equation*}
	\sum V = 0
\end{equation*}

Recorrendo aos valores da tabela 6:
\begin{equation*}
	(-13.95 \pm 0,01) - (1.169 \pm 0,01) + (15.17 \pm 0,01)
	= (+0,051 \pm 0,03)\ \text{V}
\end{equation*}

Considerando a segunda malha defenida pelo circuito constituído pelas resistências $R_2$ e $R_3$, temos que:
\begin{equation*}
	(-1.169 \pm 0,01) + (1.169 \pm 0,01)
	= (\pm 0,0 \pm 0,02)\ \text{V}
\end{equation*}

Os resultados obtidos confirmam a lei das malhas, uma vez que a soma algébrica das tensões em cada malha é aproximadamente igual a zero. Na primeira malha observa-se um pequeno desvio em relação a zero, enquanto na segunda malha o resultado é praticamente nulo. Esta diferença pode ser explicada pelas incertezas associadas às medições experimentais, como a resolução e precisão do multímetro ($\pm 0{,}01$ V), as tolerâncias dos componentes do circuito, possíveis resistências internas dos fios e da fonte, bem como pequenas variações ou ruído durante a medição. Assim, considerando estas limitações experimentais, os resultados obtidos são consistentes com a lei das malhas de Kirchoff.

\vspace{1 cm}

\subsection{Montagem 3}

\subsubsection{Verificação do príncipio da sobreposição:}

De acordo com o princípio da sobreposição (ver secção 1.1), sendo o circuito final alimentado tanto pela fonte como pela pilha, espera-se que o efeito combinado de ambas as fontes em qualquer uma das resistências seja igual à soma dos efeitos parciais provocados por cada uma das fontes individualmente, isto é num circuito com apenas uma das fontes.


Assim, arbitrariamente optou-se por realizar as medições na resistência $R_2$, pelo que para a corrente $I_2$ temos que:

\[
	I_2 = I_2^{15V} + I_2^{7.5V}
\]

\textbf{Com apenas a fonte de $15$ V ligada} (fig.~3), mediu-se:

\[
	I_2^{15V} = (13,73 \pm 0,01)\ \text{mA}
\]

\textbf{Com apenas a pilha de $7,5$ V ligada} (fig.~4), obteve-se:

\[
	I_2^{7.5V} = (0,91 \pm 0,01)\ \text{mA}
\]

Somando as contribuições individuais das duas fontes:

\[
	I_2^{15V} + I_2^{7.5V}
	= (13,73 + 0,91)\ \text{mA}
	= (14,64 \pm 0,02)\ \text{mA}
\]

\textbf{Com a fonte de $15$ V e a pilha de $7,5$ V ligadas simultaneamente} (fig.~5), obteve-se experimentalmente:

\[
	I_2 = (14,70 \pm 0,01)\ \text{mA}
\]


Comparando os resultados, verifica-se que o valor obtido experimentalmente \textbf{não} se encontra dentro da gama de valores determinada pelo valor obtido indiretamente pela soma da corrente $I_3$ nas montagens das figs. 3 e 4, pelo que o princípio da sobreposição \textbf{não foi} verificado.\\
\\De forma análoga para a tensão de $R_2$:

\[
	V_2 = V_2^{15V} + V_2^{7.5V}
\]

\textbf{Com apenas a fonte de $15$ V}:

\[
	V_2^{15V} = (13,95 \pm 0,01)\ \text{V}
\]

\textbf{Com apenas a pilha de $7,5$ V}:

\[
	V_2^{7.5V} = (0,92 \pm 0,01)\ \text{V}
\]

Logo,

\[
	V_2^{15V} + V_2^{7.5V} =
	(14,87 \pm 0,02)\ \text{V}
\]

\textbf{Com a fonte de tensão e a pilha ligadas}, mediu-se:

\[
	V_2 = (14,86 \pm 0,01)\ \text{V}
\]

Assim, dentro das incertezas experimentais das medições realizadas, os resultados obtidos para a tensão $V_2$ \textbf{são}, ao contrário da análise para a corrente $I_2$, \textbf{consistentes} com o princípio da sobreposição.


